\documentclass[11pt,a4paper]{article}
\usepackage{fontspec}
\setmainfont{Noto Serif}
\setsansfont{Noto Sans}
\usepackage{babel}
\babelprovide[main,import]{russian}
\usepackage{amsmath,amssymb,mathtools,booktabs,geometry,fancyhdr,microtype,xcolor,hyperref}
\geometry{left=21mm,right=21mm,top=21mm,bottom=21mm}
\definecolor{ink}{HTML}{173D57}
\hypersetup{colorlinks=true,urlcolor=ink,linkcolor=ink,pdftitle={BGS-F. Закон расщепления двойных уровней}}
\pagestyle{fancy}\fancyhf{}\lhead{\small BGS-F}\rhead{\small Закон расщепления двойных уровней}\cfoot{\thepage}
\setlength{\headheight}{15pt}\setlength{\parindent}{1em}\setlength{\parskip}{2.5pt}
\setlength{\emergencystretch}{3em}\tolerance=3000
\newcommand{\Tr}{\operatorname{Tr}}
\newcommand{\EE}{\mathbb E}
\newcommand{\R}{\mathbb R}
\newcommand{\Z}{\mathbb Z}
\newcommand{\C}{\mathbb C}
\newcommand{\si}{\mathrm{Sine}_2}
\newcommand{\eps}{\varepsilon}
\newcommand{\sinc}{\operatorname{sinc}}
\newcommand{\sect}[1]{\section*{\sffamily\color{ink}#1}}
\newcommand{\subt}[1]{\par\medskip\noindent\textbf{#1.}\quad}
\newcommand{\T}{\mathbb T}
\newcommand{\dist}{\operatorname{dist}}
\newcommand{\Ad}{\operatorname{Ad}}
\newcommand{\op}{\mathrm{op}}
\newcommand{\HS}{\mathrm{HS}}
\newcommand{\Zp}{\mathbb Z_7}
\newcommand{\Qp}{\mathbb Q_7}
\newcommand{\as}{\operatorname{Arcsine}}
\begin{document}
\noindent{\small\sffamily\color{ink} МАТЕМАТИЧЕСКАЯ ФИЗИКА\hfill 20 НОЯБРЯ 2026}
\begin{center}{\LARGE\sffamily Закон расщепления двойных уровней}\\[4pt]
{\large Арксинусное распределение щели на степенях семи}
\end{center}
Пусть \(N=7^k\), \(m=3N/7\) и \(U_N(t)=e^{itA_N}M_N\) для исходных операторов \(A_N,M_N\). Физическая сила \(u_N=t/N\). Выбираем равномерно один из \(m/7\) линейных двойных узлов, отвечающих характерам с \(\ell\equiv3\bmod7\); их точное определение в разделе 1. Это одна седьмая всех линейных узлов и \(6N/49\) уровней с учётом кратности.

\subt{Щель настоящего спектра}
Пусть \(t=t_N\ne0\) и
\[
                       |t_N|\sqrt N\log\log N\longrightarrow0.
\]
С вероятностью \(1-o(1)\) около выбранного узла есть ровно две собственные фазы \(U_N(t_N)\). Обозначим угловую щель между ними в радианах через \(\Delta\theta_N\). Тогда
\[
       \frac{\sqrt N}{|t_N|}\Delta\theta_N
                  \ \Longrightarrow\ \frac7{300}\,|2Y_1-Y_2-Y_3|,
              \qquad Y_i\ \text{независимы},\quad
              d\as(y)=\frac{\mathbf1_{|y|<1}}{\pi\sqrt{1-y^2}}\,dy .
\]
Предел имеет поддержку \([0,7/75]\), его второй момент \(49/30000\).
Вероятность щели не больше \(\eta\) в этом предельном законе линейна по \(\eta\) при \(\eta\downarrow0\), с конечным положительным коэффициентом.

\subt{Совместное смещение}
Для обеих фаз выведен совместный закон после деления их смещений от центра узла на \(t_N/\sqrt N\). В нём присутствует ещё одна независимая равномерная фаза, отвечающая двум непримитивным меткам кика. Её вклад общий для уровней и выпадает из щели. Предельное распределение определяется четырьмя различными квадратичными нормами меток в расщеплённых координатах.

\subt{Область результата}
Предыдущая локальная теорема давала предельную решётку с кратностями. Здесь разрешается расщепление внутри двойного узла в той же области \(u_N=o(N^{-3/2}/\log\log N)\). Модель и способ счёта фаз не изменены. Масштаб \(u_N\asymp N^{-3/2}\) и гипотеза BGS-F при фиксированном \(u_N=10^{-80}\) остаются открытыми.
\newpage
\sect{1. Выбранные двойные узлы и стационарные точки}
В расщеплённых симплектических координатах метка \((v_1,v_2)\) становится
\[
     (a,b)=\bigl(v_1+\rho v_2,\,(v_1-\rho v_2)/(-2\rho)\bigr),
                \quad \rho^2=2,\quad \rho\equiv3\bmod7 .
\]
Линейный оператор там совпадает с дилатацией \(x\mapsto\gamma^{-1}x\) с точностью до общей фазы, \(\gamma=3+2\rho\) имеет порядок \(m\). Квадраты и неквадраты первого разряда образуют две орбиты единиц. Пусть \(\omega:\mathbb F_7^\times\to\{1,\zeta_3,\zeta_3^2\}\) --- характер с \(\omega(3)=\zeta_3\), а
\[
 \mathcal L(x)=\tfrac16\log(x^6),\qquad
 \chi_{\ell,h}(x)=\omega(x\bmod7)^h e_N(\ell\mathcal L(x)),
        \quad\ell\bmod7^{k-1},\ h\bmod3 .
\]
Величина \(\mathcal L(x)\) принадлежит \(7\Zp\), в \(e_N\) берётся её остаток по модулю \(N\). Когда \(7\nmid\ell\), соответствующая фаза \(\lambda_j\) имеет кратность \(2\). Её векторы \(u_+,u_-\) равны \(m^{-1/2}\chi_{\ell,h}\) на квадратной и неквадратной орбитах соответственно. Слои меньшего периода её не содержат. Все пары \((\ell,h)\) с \(\ell\equiv3\bmod7\) дают \(3\cdot7^{k-2}=m/7\) различных узлов. Эти факты следуют из диагонализации, приведённой в части о кратностях при фиксированном кике.

Компрессия \(E_j A_N E_j\) по спектральному проектору \(E_j\) определяет производные поднятых углов при \(t=0\). Для примитивной метки в её матричном элементе фаза равна
\[
 \Phi_{\ell,v}(x)=b(x+a/2)
        +\ell\,\mathcal L((x+a)/x),\qquad
 \Phi_{\ell,v}'(x)=b-\frac{\ell a}{x(x+a)} .
\]
Стационарность по модулю \(7\) требует \(b x(x+a)=\ell a\), причём \(x,x+a\) --- единицы. При \(\ell\equiv3\) прямая проверка \(32\) меток даёт ровно шесть примитивных стационарных вкладов из квадратной орбиты. Их метки
\[
               \pm v_1,\ \pm v_2,\ \pm v_3,\qquad
        v_1=(0,1),\quad v_2=(0,8),\quad v_3=(0,-6).
\]
У каждой одна стационарная точка по модулю \(7\), \(x+a\) лежит в той же квадратной орбите, и \(2x+a\) обратимо. У \(v_1,v_2,v_3\) отношение \((x+a)/x\) равно \(4\bmod7\). Таблица проверяется простым циклом по \(x\in\{1,2,4\}\) и включена в архив. Для пары непримитивных меток \(\pm v_4,\ v_4=(0,7)\), обе координаты делятся на \(7\). После снятия множителя \(7\) стационарное уравнение равно \(b_0 x(x+a)=\ell a_0\), где \(a=7a_0,\ b=7b_0\). В квадратной орбите оно имеет один простой корень \(x\equiv2\bmod7\). Внедиагональный переход этой пары невозможен, поскольку \(a\equiv0\).

Во всех остальных классах разделение \(x=x_0+7^{\lceil k'/2\rceil}y\) обнуляет соответствующие суммы, где \(k'=k\) для примитивной метки и \(k'=k-1\) после снятия \(7\). В частности, \(E_j A_N E_j\) диагональна в базисе \(u_+,u_-\) при \(k\ge3\).
\newpage
\sect{2. Точная формула и четыре критические фазы}
Пусть \(x_i(\ell)\) --- поднятая по Гензелю стационарная точка метки \(v_i\) в квадратной орбите, \(\Phi_i(\ell)=\Phi_{\ell,v_i}(x_i(\ell))\). При примитивной метке сумма по её классу \(x\) равна \(N^{1/2}\varepsilon_i(k,h)e_N(\Phi_i)\) с \(|\varepsilon_i|=1\). Действительно, на диске \(x=x_i+7y\) гладкая \(7\)-адическая замена сводит фазу к \(\Phi_i+7^2 c_i^{\rm Hess}y^2\), \(c_i^{\rm Hess}\) обратимо. Квадратичная сумма имеет модуль \(7^{(k-2)/2}\), а число повторений равно \(7\). Замена получается из \((1+7y B(y))^{1/2}\) в разложении логарифма и биективна по каждому модулю. Это тот же гауссов расчёт, что в части о кратностях; общую форму метода см. \href{https://doi.org/10.1006/jnth.2002.2790}{B. Fisher, J. Number Theory 96 (2002), 201--224}. Для \(v_4\) исходная фаза имеет множитель \(7\); после снятия модуля \(7^{k-1}\) добавляются семь повторений. Её ответ имеет модуль \((7N)^{1/2}\). Факторы \(\varepsilon_i(k,h)\) зависят только от \(k,h\) и исходного класса \(\ell\bmod7=3\). Эти множители образуются из нормированных квадратичных сумм и конечного характера отношения \((x+a)/x\bmod7\).

Вейлев оператор отрицательной метки сопряжён к положительному. Чётность \(x\mapsto-x\) меняет \(u_+,u_-\) и \(T(v)\) на \(T(-v)\), поскольку \(\chi_{\ell,h}(-1)=1\). Поэтому, обозначая через \(c_i=c_{v_i}\) коэффициенты исходного кика и \(Z_{i,N}=\varepsilon_i(k,h)e_N(\Phi_i)\), получаем два диагональных матричных элемента после умножения на \(\sqrt N\):
\begin{equation}
 \begin{aligned}
 D_{+,N}&=\frac{14}{3}\Re\left(\sum_{i=1}^3 c_i Z_{i,N}
                                           +\sqrt7\,c_4Z_{4,N}\right),\\
 D_{-,N}&=\frac{14}{3}\Re\left(\sum_{i=1}^3 \bar c_i Z_{i,N}
                                           +\sqrt7\,c_4Z_{4,N}\right).
 \end{aligned}                                                   \tag{1}
\end{equation}
Здесь \((c_1,c_2,c_3,c_4)=((200-20i),(-100+10i),(-100+10i),1800)/4000\). Нормировка \(14/3\) учитывает \(m^{-1}\), модуль \(N^{1/2}\) и добавление сопряжённого вклада.

\subt{Аналитическая зависимость от \(\ell\)}
Функции \(x_i(\ell)\) аналитичны на \(\ell\in3+7\Zp\). Это следует из простого корня уравнения после снятия общего множителя \(7\) у \(v_4\); квадратный корень единицы вычисляется биномиальным рядом на одном классе. Положим \(q_i=a_i b_i=\rho\,v_{i,2}^2/2\) и \(z_i=x_i/a_i\). Даже для \(v_4\), где \(z_i\in\Qp\) с отрицательным порядком, верно \(q_i z_i(z_i+1)=\ell\). Дифференцируя критическую фазу по \(\ell\) и используя стационарность по \(x\), получаем
\begin{equation}
          \Phi_i'(\ell)=\mathcal L((z_i+1)/z_i),\qquad
          \ell\,\Phi_i''(\ell)=-\frac1{2z_i+1},\qquad
                   (2z_i+1)^2=1+4\ell/q_i .                     \tag{2}
\end{equation}
Отношение в логарифме всегда единица. Неравенство \(q_i\ne q_j\) при \(i\ne j\) следует из \(v_{i,2}^2\in\{1,64,36,49\}\).
\newpage
\sect{3. Совместная равномерность и закон щели}
Выберем \(\ell\equiv3\bmod7\) равномерно по модулю \(7^{k-1}\). Докажем \((Z_{1,N},\ldots,Z_{4,N})\Rightarrow (Z_1,\ldots,Z_4)\), где \(Z_i\) независимы и равномерны на единичной окружности. Для \(\mathbf n\in\Z^4\setminus\{0\}\) положим \(F=\sum_i n_i\Phi_i\). Из (2) функции \(g_i=-(2z_i+1)^{-1}\) удовлетворяют
\[
 g_i^{(r)}(\ell)=g_i(\ell)
       \prod_{a=0}^{r-1}(-\tfrac12-a)\left(\frac4{q_i+4\ell}\right)^r .
\]
Матрица этих значений с \(0\le r\le3\) невырождена. После деления строк и столбцов на ненулевые множители это матрица Вандермонда по четырём различным \(4/(q_i+4\ell)\). Все \(q_i+4\ell\) ненулевые в выбранном классе. Поэтому ни в одной точке одновременно не исчезают первые четыре производные порядков \(0,1,2,3\) функции \(\ell F''\). Любой нуль \(F'\) имеет порядок не больше четырёх. Нули изолированы, и на компактном диске их конечное число. В окрестности каждого нуля аналитическое разложение даёт
\[
               \mathbb P_{3+7\Zp}\bigl(\nu_7(F')\ge s\bigr)
                              =O_{\mathbf n}(7^{-s/4}).
\]
На дополнении к этим окрестностям порядок ограничен по компактности.

Так как \(\Phi_i'\) кратна \(7\), значение \(e_N(F(\ell))\) зависит только от \(\ell\bmod7^{k-1}\). Для точного сокращения выбираем \(\ell\) по модулю \(7^k\) и раскладываем \(\ell=\ell_*+7^{\lceil k/2\rceil}y\). На любой такой подокрестности корень \(x_i\) имеет интегральный ряд по приращению. Для логарифма относительных приращений \(\log(1+g)=g+O_7(7^{2s})\), если \(g\in7^s\Zp\), \(s\ge1\), поскольку \(rs-\nu_7(r)\ge2s\) при \(r\ge2\). Поэтому \(F(\ell)\equiv F(\ell_*)+7^{\lceil k/2\rceil}yF'(\ell_*)\bmod7^k\). Среднее по \(y\) равно нулю, если \(F'(\ell_*)\not\equiv0\bmod7^{\lfloor k/2\rfloor}\). Производная стабильна по этому модулю на подокрестности, так что оставшаяся доля \(O_{\mathbf n}(7^{-k/8})\). Домножение на \(\varepsilon_i(k,h)\) не меняет границу. По критерию Вейля это доказывает совместную равномерность для каждого \(h\) и для их смеси.

Из (1) предельная пара \((D_+,D_-)\) получается заменой \(Z_{i,N}\) на независимые \(Z_i\). Её абсолютная разность равна \(\tfrac7{300}|2Y_1-Y_2-Y_3|\), \(Y_i=\Im Z_i\sim\as\), так как \(c_4\) вещественен. Поддержка и второй момент следуют из \(|Y_i|\le1\), \(\mathbb EY_i=0\), \(\mathbb EY_i^2=1/2\). Для \(S=2Y_1-Y_2-Y_3\) характеристическая функция \(J_0(2s)J_0(s)^2\) интегрируема. Здесь \(J_0(s)=(2\pi)^{-1}\int_0^{2\pi} e^{is\sin x}\,dx=O((1+|s|)^{-1/2})\) по двум стационарным точкам и интегрированию по частям вне их окрестностей. Поэтому \(S\) имеет непрерывную плотность \(f_S\). Она строго положительна в нуле, например по свёртке трёх положительных плотностей на малых интервалах около нуля. Отсюда
\[
       \mathbb P\bigl((7/300)|S|\le \eta\bigr)
           = \frac{600}{7}\, f_S(0)\,\eta+o(\eta).
\]
В частности, предельная разность не имеет атома в нуле.
\newpage
\sect{4. Устранение дальнего блока с точностью до масштаба щели}
Положим \(\epsilon=|t|\sqrt N\), \(L=\lfloor k^{16}\rfloor\), \(P=E_j\), \(Q=\sum_{1\le |i-j|_m\le L}E_i\), \(F=I-P-Q\). В этом и следующем разделе вероятность берётся по равномерному выбору \(j\) из выделенного класса; вся оценка среднего по всем \(j\) переходит на него с множителем \(7\). Символ \(O_{\mathbb P}\) обозначает ограниченность по вероятности при \(k\to\infty\), равномерно по \(z\) в области, если это указано. Пусть
\[
   \Omega_j=\{z:|z-\lambda_j|\le1/m\},\qquad
   D_z=\sqrt m\,P+ (M_N-zI)^{-1/2}\big|_{Q+F}.
\]
Корень на \(Q+F\) выбирается аналитически в окрестности диска. После умножения \(U_N(t)-zI\) на \(D_z\) с обеих сторон невозмущённая часть равна \(m(\lambda_j-z)P+Q+F\). В этой области положительные резольвентные веса как в части о знаковом убывании с радиусом \(1\) дают \(|D_z|\le\sqrt m Z\). Обозначим \(Z_{\rm n}=(P+Q)Z\), \(Z_{\rm f}=FZ\), \(B^{(1)}=itD_z A_N M_N D_z\).

Из четвёртых оценок
\[
 \EE_j\|tm Z_{\rm f} A_N Z_{\rm f}\|_{S_4}^4
          \le C\epsilon^4 k^3/L,\qquad
 \EE_j\|tm Z A_N Z\|_{S_4}^4\le C\epsilon^4 k^2
\]
следует, что вне доли \(o(1)\) линейный дальний блок \(\le\epsilon k^{-2}\) по норме, а обе смешанные нормы \(\le\epsilon k^{3/4}\). Действительно, Марков даёт \(O(k^{11}/L)\) и \(O(k^{-1})\) соответственно. Это также выполняется равномерно по \(z\) через диагональную мажоранту. Нелинейный остаток \(e^{itA_N}-I-itA_N\) после умножения имеет порядки \(O(\epsilon^2)\), \(O(\epsilon^2/L)\), \(O(\epsilon^2/\sqrt L)\) в ближнем, дальнем и смешанном блоках соответственно.

После устранения дальнего блока \(F\) его дополнение Шура в ближней части можно заменить на произведение двух линейных смешанных блоков с ошибкой \(o_{\mathbb P}(\epsilon^2)\). Замена дальнего обратного на \(I\) стоит \(O(\epsilon^3 k^{-1/2})\), замена смешанных блоков на линейные --- \(O(\epsilon^3 k^{3/4}/\sqrt L+\epsilon^4/L)\). В линейном произведении внутреннюю дальнюю резольвенту можно заменить на \(mF_j\), где
\[
                F_j=\frac1m F M_N (M_N-\lambda_j I)^{-1}\big|_F.
\]
Это стоит \(O(\epsilon^2k^{3/2}/L)\) равномерно на \(\Omega_j\) по тождеству резольвент и относительной ошибке \(O(1/L)\) на каждой дальней фазе. Для центрального произведения применяем усиленную оценку знаковой резольвенты
\[
        \EE_j\|t^2m^2Z_{\rm n} A_N F_j A_N Z_{\rm n}\|_{\HS}^2
                                 \le C\epsilon^4 .
\]
Итак, после устранения \(F\) ближняя матрица равна
\begin{equation}
      m(\lambda_j-z)P+Q+(P+Q)B^{(1)}(P+Q)
                              +O_{\mathbb P}(\epsilon^2),
                                                               \tag{3}
\end{equation}
равномерно по диску. Эта относительная оценка важна и при очень малых ненулевых \(t\), так как весь ведущий сдвиг имеет нормированный порядок \(\epsilon\).
\newpage
\sect{5. Оставшийся блок и истинные фазы}
Для всех разностей \(d\) в линейном спектре введём \(W_d=\sum_i\|E_i A_N E_{i+d}\|_{\HS}^2\). Как было доказано через проекции на орбиты вейлевых меток, \(W_d\le C N/m\). По двум весам в ближнем блоке и одному весу в переходе из центра,
\begin{align*}
 \EE_j\|tm Z_{\rm n} A_N Z_{\rm n}\|_{\HS}^2
       &\le C\, t^2N\,(1+\log(2L))^2,\\
 \EE_j\|tm P A_N QZ_{\rm n}\|_{\HS}^2
       &\le C\, t^2 m\sum_{1\le |d|_m\le L}w_d W_d
        \le C\,t^2N\,(1+\log(2L)).
\end{align*}
Вторая оценка следует раскрытием квадрата и усреднением по \(j\). Здесь у центрального проектора вес равен \(1\), поэтому остаётся одна гармоническая сумма. Эрмитовость \(A_N\) даёт ту же границу для сопряжённого перехода. Диагональные факторы \(D_z\) и \(M_N\) переносят эти оценки на линейную часть (3) равномерно по \(\Omega_j\).

Если \(\epsilon\log k\to0\), блок \(Q\) в (3) отличается от единицы на \(o_{\mathbb P}(1)\), поэтому обратим с вероятностью \(1-o(1)\) и равномерно по \(z\). Каждая из двух норм перехода \(P\leftrightarrow Q\) есть \(O_{\mathbb P}(\epsilon\sqrt{1+\log k})\). Следовательно, дополнение Шура на двумерном \(P\) равно
\begin{equation}
 m(\lambda_j-z)I_P+itm\lambda_j E_j A_N E_j
              +O_{\mathbb P}(\epsilon^2(1+\log k)).
                                                               \tag{4}
\end{equation}
Остаток в (4) есть \(o_{\mathbb P}(\epsilon)\). Обе степени исключения \(F,Q\) корректны при каждом \(z\) на всём диске, а их определители не обращаются в нуль. На \(P\) умножение на \(D_z\) есть фиксированный скаляр \(\sqrt m\), так что нули определителя (4) совпадают с корнями исходной матрицы в \(\Omega_j\).

\subt{Определение числа фаз и сравнение смещений}
На внешнем контуре \(|z-\lambda_j|=1/m\) первый член (4) имеет модуль \(1\); второй \(O(\epsilon)\) по точной формуле (1). Принцип аргумента через линейное добавление остальных членов показывает, что внутри ровно две собственные фазы с вероятностью \(1-o(1)\). В масштабе
\[
                       x=\frac{\sqrt N(z-\lambda_j)}{i t\lambda_j}
\]
после деления на \(itm\lambda_j/\sqrt N\) матрица (4) становится \(H_j-xI+o_{\mathbb P}(1)\), где \(H_j=\sqrt N E_j A_N E_j\) эрмитова. Её собственные числа \(D_{\pm,N}\) равномерно ограничены. На двух окружностях фиксированного малого радиуса \(\delta\) в \(x\)-переменной вокруг \(D_{\pm,N}\), если разность больше \(2\delta\), обратная норма \((H_j-xI)^{-1}\le1/\delta\). При таком выборе внешний и обе внутренние окружности лежат в \(\Omega_j\) для больших \(k\). Принцип аргумента даёт по одному собственному числу в каждом малом круге, если остаток меньше \(\delta\). Вероятность отказа из-за малой разности можно сделать сколь угодно малой, поскольку предельная разность не имеет атома в нуле. Это доказывает соответствие упорядоченных смещений собственным числам \(H_j\) с ошибкой \(o_{\mathbb P}(1)\) после нормировки. На единичной окружности переход от \(x\) к \(\sqrt N\arg(z/\lambda_j)/t\) меняет его на \(o_{\mathbb P}(1)\), по разложению \(e^{i\theta}-1=i\theta+O(\theta^2)\) и ограниченности \(x\) по вероятности.

Из раздела 3 следует утверждение о настоящих фазах и их щели. Масштаб \(|t|/\sqrt N=|u_N|\sqrt N\) меньше средней угловой щели \(2\pi/N\) при указанном условии. Это распределение разрешает кратности прежнего локального предела внутри выбранных кластеров.
\newpage
\sect{6. Проверки и точный статус}
В приложенном коде отдельно проверены исходные коэффициенты кика, остатки стационарных корней, отсутствие иных примитивных переходов при \(\ell\equiv3\bmod7\) и сохранение фазы при подъёме \(\ell\) по модулю \(7^{k-1}\). Все эти операции выполняются целыми числами. Для проверки (1) вычисляются прямые матричные элементы в расщеплённых координатах при \(N=343,2401\), с независимой суммой всех \(32\) меток, включая непримитивные. При \(N=343\) также диагонализуется \(e^{itA_N}M_N\) для вспомогательного \(t=2\cdot10^{-6}\) и сверяются обе нормированные угловые фазы около каждого выбранного тестового узла. При выполнении проверки максимальная ошибка была меньше \(6\cdot10^{-6}\). Численный тест проверяет конвенции и формулу. Теорема опирается на точный расчёт и асимптотические оценки.

\subt{Что остаётся}
Улучшенная оценка дальнего знакового блока имеет порядок \(\epsilon^2\). Для выявления истинных фаз после отсечения дальних узлов используется оценка ближнего блока \(O_{\mathbb P}(\epsilon\log k)\), и она требует \(\epsilon\log k\to0\). Равномерная асимптотика при \(\epsilon\asymp1\), то есть для \(u_N\asymp N^{-3/2}\), здесь не выводится. При фиксированной исходной силе \(\epsilon=N^{3/2}/10^{80}\to\infty\), поэтому основная гипотеза BGS-F остаётся открытой.
\end{document}
